\documentclass[11pt]{article}

\usepackage[a4paper,margin=1in]{geometry}
\usepackage{amsmath,amssymb,amsthm,mathtools}
\usepackage[T1]{fontenc}
\usepackage{lmodern}
\usepackage{microtype}

% Fallback when hyperref is not loaded.
\providecommand{\texorpdfstring}[2]{#1}

\newtheorem{theorem}{Theorem}[section]
\newtheorem{lemma}[theorem]{Lemma}
\newtheorem{proposition}[theorem]{Proposition}
\newtheorem{corollary}[theorem]{Corollary}
\newtheorem{remark}[theorem]{Remark}
\newtheorem{conjecture}[theorem]{Conjecture}
\newtheorem{definition}[theorem]{Definition}

\newcommand{\C}{\mathbb C}
\newcommand{\R}{\mathbb R}
\newcommand{\eps}{\varepsilon}

\title{A Dual-Tail Local Reduction for the Real-Slice Problem\\
in the Clean OBC Hatano--Nelson Family}
\author{}
\date{}

\begin{document}
\maketitle

\section{Setup}

Fix \(0<a<b\), and let
\[
A_n=\operatorname{tridiag}(a,0,b)\in\R^{n\times n}.
\]
For \(z=x+iy\in\C\), define
\[
B(z)\coloneqq zI-A_n,
\qquad
M_n(x,y)\coloneqq B(z)^*B(z),
\qquad
\lambda_*(x,y)\coloneqq \lambda_{\min}(M_n(x,y)).
\]
At a boundary point of the pseudospectrum,
\[
\Lambda=\lambda_*(x,y),
\]
one has
\[
M_n(x,y)-\Lambda I_n \succeq 0,
\qquad
\det(M_n(x,y)-\Lambda I_n)=0.
\]

We write
\[
q\coloneqq b-a>0,
\qquad
\eta\coloneqq ab,
\qquad
c\coloneqq (a+b)x-iqy.
\]

The connectedness-equivalence problem is to prove that for every fixed \(x\in\R\) and every
\(\eps>0\), the vertical section
\[
\{y\in\R:\ \lambda_*(x,y)<\eps^2\}
\]
is either empty or an interval containing \(0\). It suffices to prove that, at every simple
boundary point with \(y>0\),
\[
\partial_y \lambda_*(x,y)>0.
\]

\section{Block form and dual tails}

We first record the even case explicitly. The odd case has the same interior structure and only
a modified terminal seed.

\subsection*{Even case \(n=2m\)}

Set
\[
w_k\coloneqq \binom{u_{2k-1}}{u_{2k}},
\qquad
1\le k\le m.
\]
Then
\[
\bigl(M_{2m}(x,y)-\Lambda I_{2m}\bigr)u=0
\]
is equivalent to the \(2\times2\) block Jacobi system
\[
B_L w_1 + A w_2 = 0,
\]
\[
A^* w_{k-1}+Bw_k+Aw_{k+1}=0,
\qquad
2\le k\le m-1,
\]
\[
A^*w_{m-1}+B_R w_m=0,
\]
where
\[
A=
\begin{pmatrix}
\eta & 0\\
-c & \eta
\end{pmatrix},
\]
\[
B_L=
\begin{pmatrix}
x^2+y^2+a^2-\Lambda & -c\\
-\bar c & x^2+y^2+a^2+b^2-\Lambda
\end{pmatrix},
\]
\[
B=
\begin{pmatrix}
x^2+y^2+a^2+b^2-\Lambda & -c\\
-\bar c & x^2+y^2+a^2+b^2-\Lambda
\end{pmatrix},
\]
\[
B_R=
\begin{pmatrix}
x^2+y^2+a^2+b^2-\Lambda & -c\\
-\bar c & x^2+y^2+b^2-\Lambda
\end{pmatrix}.
\]

\begin{definition}[Right and left tail Weyl matrices]
In the even case, define recursively the right tails
\[
\Sigma_m\coloneqq B_R,
\qquad
\Sigma_k\coloneqq B_k-A\Sigma_{k+1}^{-1}A^*,
\qquad
k=m-1,\dots,1,
\]
where
\[
B_1\coloneqq B_L,
\qquad
B_k\coloneqq B
\quad (2\le k\le m-1),
\]
and the left tails
\[
\Gamma_1\coloneqq B_L,
\qquad
\Gamma_k\coloneqq B_k-A^*\Gamma_{k-1}^{-1}A,
\qquad
k=2,\dots,m,
\]
where
\[
B_k\coloneqq B
\quad (2\le k\le m-1),
\qquad
B_m\coloneqq B_R.
\]
\end{definition}

\begin{proposition}[Tail action on the block nullchain]\label{prop:tails}
Let \(w=(w_1,\dots,w_m)\) be a nonzero block solution of the boundary equation
\[
(M_{2m}(x,y)-\Lambda I_{2m})u=0.
\]
Then
\[
\Sigma_k w_k=-A^*w_{k-1},
\qquad
\Gamma_k w_k=-Aw_{k+1},
\]
for every admissible \(k\), with the endpoint conventions
\[
A^*w_0=0,
\qquad
Aw_{m+1}=0.
\]
\end{proposition}

\begin{proof}
This is the standard Schur-complement interpretation of the right and left tails.
\end{proof}

\begin{definition}[Local effective block]
For \(1\le k\le m\), define
\[
K_k\coloneqq \Sigma_k+\Gamma_k-B_k,
\]
where \(B_1=B_L\), \(B_m=B_R\), and \(B_k=B\) for \(2\le k\le m-1\).
\end{definition}

\begin{proposition}[Exact local invariant]\label{prop:K}
At a boundary point \(\Lambda=\lambda_*(x,y)\), one has
\[
K_k\succeq 0,
\qquad
\operatorname{rank}K_k=1,
\qquad
K_k w_k=0
\]
for every \(k\). Therefore
\[
K_k=
\begin{pmatrix}
A_k & \zeta_k+i t_k\\
\zeta_k-i t_k & D_k
\end{pmatrix},
\]
with
\[
A_k>0,\qquad D_k>0,\qquad t_k>0,
\]
and
\[
A_kD_k=\zeta_k^2+t_k^2.
\]
If
\[
r_k\coloneqq \frac{w_{k,2}}{w_{k,1}},
\]
then
\[
r_k\in\mathbb H:=\{z\in\C:\Im z>0\},
\]
and the within-block current
\[
j_k\coloneqq \Im(\overline{w_{k,1}}\,w_{k,2})
\]
satisfies
\[
j_k=\frac{t_k}{D_k}|w_{k,1}|^2=\frac{t_k}{A_k}|w_{k,2}|^2>0.
\]
\end{proposition}

\begin{proof}
The matrix \(K_k\) is the Schur complement obtained by eliminating all blocks except block \(k\)
from the positive semidefinite matrix \(M_{2m}(x,y)-\Lambda I\). Hence \(K_k\succeq 0\). Since the
global kernel is one-dimensional, the local Schur complement also has one-dimensional kernel, so
\(\operatorname{rank}K_k=1\). The rest is immediate from the explicit \(2\times2\) form.
\end{proof}

\section{Cross-block currents and local mirror criteria}

Define the cross-block currents
\[
h_k\coloneqq \Im(\overline{w_{k,2}}\,w_{k+1,1}),
\qquad
1\le k\le m-1.
\]

Write the right tail in coordinates
\[
\Sigma_{k+1}=
\begin{pmatrix}
u_{k+1}^R & \beta_{k+1}^R\\
\overline{\beta_{k+1}^R} & v_{k+1}^R
\end{pmatrix},
\qquad
\Delta_{k+1}^R\coloneqq u_{k+1}^R v_{k+1}^R-|\beta_{k+1}^R|^2,
\qquad
\tau_{k+1}^R\coloneqq \Im \beta_{k+1}^R,
\]
and the left tail in coordinates
\[
\Gamma_{k-1}=
\begin{pmatrix}
U_{k-1}^L & \gamma_{k-1}^L\\
\overline{\gamma_{k-1}^L} & V_{k-1}^L
\end{pmatrix},
\qquad
\Delta_{k-1}^L\coloneqq U_{k-1}^L V_{k-1}^L-|\gamma_{k-1}^L|^2,
\qquad
\tau_{k-1}^L\coloneqq \Im \gamma_{k-1}^L.
\]

\begin{proposition}[Exact formulas for adjacent cross-currents]\label{prop:h}
For every interior block \(2\le k\le m-1\), one has
\[
h_k
=
\frac{
\eta v_{k+1}^R\,j_k
+
(\eta\tau_{k+1}^R+qy\,v_{k+1}^R)\,|w_{k,2}|^2
}{
\Delta_{k+1}^R
},
\]
and
\[
h_{k-1}
=
\frac{
\eta U_{k-1}^L\,j_k
+
(\eta\tau_{k-1}^L+qy\,U_{k-1}^L)\,|w_{k,1}|^2
}{
\Delta_{k-1}^L
}.
\]
In particular,
\[
h_k>0,
\qquad
h_{k-1}>0.
\]
\end{proposition}

\begin{proof}
These follow by substituting the tail identities
\[
\Sigma_{k+1}w_{k+1}=-A^*w_k,
\qquad
\Gamma_{k-1}w_{k-1}=-Aw_k,
\]
and taking imaginary parts of the appropriate scalar products.
\end{proof}

Now define the mirror-pair quantities
\[
p_j\coloneqq -\Im(u_j u_{n+1-j}),
\qquad 1\le j\le n.
\]

\begin{proposition}[Local criteria for mirror-pair positivity]\label{prop:Psi}
For every interior block \(2\le k\le m-1\), one has
\[
s\,p_{2k-1}
=
y|w_{k,1}|^2+a h_{k-1}-b j_k,
\]
\[
s\,p_{2k}
=
y|w_{k,2}|^2+a j_k-b h_k,
\]
where
\[
s=\sigma_{\min}(x+iy-A_n).
\]
Using Proposition~\ref{prop:h} and the identity
\[
j_k=\frac{t_k}{D_k}|w_{k,1}|^2=\frac{t_k}{A_k}|w_{k,2}|^2,
\]
this becomes
\[
p_{2k-1}>0
\quad\Longleftrightarrow\quad
\Psi_k^- >0,
\]
\[
p_{2k}>0
\quad\Longleftrightarrow\quad
\Psi_k^+ >0,
\]
where
\[
\Psi_k^-
=
y
+
\frac{a(\eta\tau_{k-1}^L+qy\,U_{k-1}^L)}{\Delta_{k-1}^L}
+
\left(
\frac{a\eta U_{k-1}^L}{\Delta_{k-1}^L}-b
\right)\frac{t_k}{D_k},
\]
and
\[
\Psi_k^+
=
y
-
\frac{b(\eta\tau_{k+1}^R+qy\,v_{k+1}^R)}{\Delta_{k+1}^R}
+
\left(
a-\frac{b\eta v_{k+1}^R}{\Delta_{k+1}^R}
\right)\frac{t_k}{A_k}.
\]
\end{proposition}

\begin{remark}
Thus the mirror-pair sign problem has been reduced to a local \(2\times2\) inequality at each
block:
\[
\Psi_k^- >0,\qquad \Psi_k^+>0.
\]
No global prefix sum remains in this formulation.
\end{remark}

\section{A global Takagi identity}

We now state the cleanest global scalar reduction.

\begin{proposition}[Takagi form and exact current identities]\label{prop:Takagi}
Let
\[
B(z)\coloneqq zI-A_n,
\qquad
z=x+iy,
\qquad
y>0.
\]
Choose a right singular vector \(u=(u_1,\dots,u_n)^T\neq 0\) in Takagi form:
\[
B(z)u=s\,J\overline{u},
\qquad
s=\sigma_{\min}(B(z)),
\]
where \(J\) is the reversal matrix.

Define
\[
I_j\coloneqq \Im(\overline{u_j}u_{j+1}),
\qquad
1\le j\le n-1,
\qquad
I_0=I_n=0,
\]
and
\[
p_j\coloneqq -\Im(u_j u_{n+1-j}),
\qquad
1\le j\le n.
\]
Then for every \(j=1,\dots,n\),
\[
s\,p_j
=
y|u_j|^2+a I_{j-1}-b I_j.
\]
If moreover \(\|u\|=1\), then
\[
s\sum_{j=1}^n p_j
=
y-q\sum_{j=1}^{n-1} I_j.
\]
Consequently,
\[
\partial_y \lambda_*(x,y)
=
2y-2q\sum_{j=1}^{n-1} I_j
=
2s\sum_{j=1}^n p_j.
\]
Finally,
\[
p_1=p_n,
\qquad
s\,p_n=y|u_n|^2+a I_{n-1}>0,
\]
so the endpoint mirror pair is always positive.
\end{proposition}

\begin{proof}
The Takagi equation is
\[
(x+iy)u_j-a u_{j-1}-b u_{j+1}=s\,\overline{u_{n+1-j}},
\qquad
u_0=u_{n+1}=0.
\]
Multiplying by \(\overline{u_j}\) and taking imaginary parts gives
\[
s\,p_j=y|u_j|^2+a I_{j-1}-b I_j.
\]
Summing over \(j\) gives
\[
s\sum_{j=1}^n p_j
=
y\sum_{j=1}^n |u_j|^2
+a\sum_{j=1}^n I_{j-1}
-b\sum_{j=1}^n I_j
=
y-q\sum_{j=1}^{n-1} I_j.
\]
For normalized \(u\), the Hellmann--Feynman formula gives
\[
\partial_y \lambda_*(x,y)=2y-2q\sum_{j=1}^{n-1}I_j,
\]
hence
\[
\partial_y \lambda_*(x,y)=2s\sum_{j=1}^n p_j.
\]
Finally, \(p_1=p_n\) is immediate from the definition, and at \(j=n\) one has
\[
s\,p_n=y|u_n|^2+a I_{n-1}>0.
\]
\end{proof}

\section{Connectedness equivalence: exact reduction and remaining gap}

\begin{lemma}[Topological vertical-slice criterion]\label{lem:topological}
Let \(\Omega\subset\C\) be open. If every vertical slice
\[
\{y\in\R:\ x+iy\in\Omega\}
\]
is either empty or an interval containing \(0\), then
\[
\Omega\ \text{is connected}
\quad\Longleftrightarrow\quad
\Omega\cap\R\ \text{is connected}.
\]
\end{lemma}

\begin{conjecture}[Mirror-pair negativity]\label{conj:mirror}
At every simple boundary point \(\Lambda=\lambda_*(x,y)\) with \(y>0\), one has
\[
p_j>0
\qquad (1\le j\le n).
\]
Equivalently, by Proposition~\ref{prop:Psi}, for every interior block \(k\),
\[
\Psi_k^- >0,
\qquad
\Psi_k^+ >0.
\]
\end{conjecture}

\begin{theorem}[Conditional completion of the real-slice criterion]\label{thm:conditional}
Assume Conjecture~\ref{conj:mirror}. Then for every \(n\ge 2\) and every \(\eps>0\),
\[
\sigma_\eps(A_n)\ \text{is connected}
\quad\Longleftrightarrow\quad
\sigma_\eps(A_n)\cap\R\ \text{is connected}.
\]
\end{theorem}

\begin{proof}
Fix \(x\in\R\) and set
\[
F(y)\coloneqq \lambda_*(x,y)-\eps^2.
\]
Then \(F\) is even and continuous, and \(F(y)\to+\infty\) as \(|y|\to\infty\).
At a simple positive root \(F(y_0)=0\), Proposition~\ref{prop:Takagi} gives
\[
F'(y_0)=\partial_y \lambda_*(x,y_0)=2s\sum_{j=1}^n p_j.
\]
Under Conjecture~\ref{conj:mirror}, this is strictly positive. Hence every positive root is
simple and crossed from negative to positive. Therefore \(\{y\ge 0:\ F(y)<0\}\) is either
empty or an interval of the form \([0,y_*)\). By evenness,
\[
\{y\in\R:\ F(y)<0\}
\]
is either empty or an interval containing \(0\). Thus every vertical slice of
\(\sigma_\eps(A_n)\) is either empty or an interval containing \(0\). The result follows from
Lemma~\ref{lem:topological}.
\end{proof}

\section{Current status}

\begin{remark}
The exact progress is the following.

\smallskip
\noindent
\textup{(1)} The finite-\(n\) problem has been reduced to the \(2\times2\) parity Weyl matrices
\(\Sigma_k\) and \(\Gamma_k\), and then to the local effective block
\[
K_k=\Sigma_k+\Gamma_k-B_k.
\]

\smallskip
\noindent
\textup{(2)} The local matrix \(K_k\) is positive semidefinite of rank \(1\), annihilates the
block solution \(w_k\), and determines the within-block current \(j_k\).

\smallskip
\noindent
\textup{(3)} The adjacent cross-currents \(h_{k-1}\) and \(h_k\) are explicit affine functions of
the local current \(j_k\) and the left/right tails.

\smallskip
\noindent
\textup{(4)} The mirror-pair sign problem has therefore been reduced to the local inequalities
\[
\Psi_k^- >0,\qquad \Psi_k^+>0.
\]

\smallskip
\noindent
\textup{(5)} Globally, the boundary slope is exactly
\[
\partial_y\lambda_*(x,y)=2s\sum_{j=1}^n p_j,
\]
so the connectedness-equivalence theorem is reduced to showing positivity of the mirror-pair
sum, and the stronger pointwise positivity conjecture above would suffice.

\smallskip
\noindent
\textup{(6)} The endpoint pair is already proved:
\[
p_1=p_n>0.
\]

\smallskip
Thus the remaining gap is now local and finite-dimensional: one must prove
\[
\Psi_k^- >0,\qquad \Psi_k^+>0
\]
for all blocks \(k\).
\end{remark}

\section{Direct calculation of the local criteria \texorpdfstring{$\Psi_k^\pm$}{Psi}}

We work at a boundary point
\[
\Lambda=\lambda_{\min}(M_n(x,y)),
\qquad
y>0.
\]
Write
\[
q\coloneqq b-a>0,
\qquad
\eta\coloneqq ab,
\qquad
c\coloneqq (a+b)x-iqy.
\]
In the even case \(n=2m\), let
\[
w_k\coloneqq \binom{u_{2k-1}}{u_{2k}},
\qquad
1\le k\le m,
\]
and let \(\Sigma_k\) and \(\Gamma_k\) denote the right and left \(2\times2\) tail Weyl matrices,
so that
\[
\Sigma_k w_k=-A^*w_{k-1},
\qquad
\Gamma_k w_k=-Aw_{k+1},
\]
with
\[
A=
\begin{pmatrix}
\eta & 0\\
-c & \eta
\end{pmatrix}.
\]

For interior blocks \(2\le k\le m-1\), write
\[
\Sigma_{k+1}=
\begin{pmatrix}
u_{k+1}^R & \beta_{k+1}^R\\
\overline{\beta_{k+1}^R} & v_{k+1}^R
\end{pmatrix},
\qquad
\beta_{k+1}^R=b_{k+1}^R+i\tau_{k+1}^R,
\qquad
\Delta_{k+1}^R\coloneqq
u_{k+1}^R v_{k+1}^R-\bigl|\beta_{k+1}^R\bigr|^2,
\]
and
\[
\Gamma_{k-1}=
\begin{pmatrix}
U_{k-1}^L & \gamma_{k-1}^L\\
\overline{\gamma_{k-1}^L} & V_{k-1}^L
\end{pmatrix},
\qquad
\gamma_{k-1}^L=g_{k-1}^L+i\tau_{k-1}^L,
\qquad
\Delta_{k-1}^L\coloneqq
U_{k-1}^L V_{k-1}^L-\bigl|\gamma_{k-1}^L\bigr|^2.
\]

Also define the local effective block
\[
K_k\coloneqq \Sigma_k+\Gamma_k-B,
\]
where
\[
B=
\begin{pmatrix}
x^2+y^2+a^2+b^2-\Lambda & -c\\
-\bar c & x^2+y^2+a^2+b^2-\Lambda
\end{pmatrix}
\]
for interior blocks. At the boundary, \(K_k\succeq 0\), \(\operatorname{rank}K_k=1\), and
\[
K_k=
\begin{pmatrix}
A_k & \zeta_k+i t_k\\
\zeta_k-i t_k & D_k
\end{pmatrix},
\qquad
A_k>0,\quad D_k>0,\quad t_k>0.
\]
If
\[
j_k\coloneqq \Im(\overline{w_{k,1}}\,w_{k,2}),
\]
then
\[
j_k=\frac{t_k}{D_k}|w_{k,1}|^2=\frac{t_k}{A_k}|w_{k,2}|^2.
\]

\begin{proposition}[Direct formulas for the adjacent cross-currents]
For every interior block \(2\le k\le m-1\), define
\[
h_k\coloneqq \Im(\overline{w_{k,2}}\,w_{k+1,1}),
\qquad
h_{k-1}\coloneqq \Im(\overline{w_{k-1,2}}\,w_{k,1}).
\]
Then
\[
h_k
=
\frac{
\eta v_{k+1}^R\,j_k
+
(\eta\tau_{k+1}^R+qy\,v_{k+1}^R)\,|w_{k,2}|^2
}{
\Delta_{k+1}^R
},
\]
and
\[
h_{k-1}
=
\frac{
\eta U_{k-1}^L\,j_k
+
(\eta\tau_{k-1}^L+qy\,U_{k-1}^L)\,|w_{k,1}|^2
}{
\Delta_{k-1}^L
}.
\]
In particular,
\[
h_k>0,
\qquad
h_{k-1}>0.
\]
\end{proposition}

\begin{proof}
From the right-tail identity
\[
w_{k+1}=-\Sigma_{k+1}^{-1}A^*w_k,
\]
with
\[
\Sigma_{k+1}^{-1}
=
\frac1{\Delta_{k+1}^R}
\begin{pmatrix}
v_{k+1}^R & -\beta_{k+1}^R\\
-\overline{\beta_{k+1}^R} & u_{k+1}^R
\end{pmatrix},
\qquad
A^*=
\begin{pmatrix}
\eta & -\bar c\\
0 & \eta
\end{pmatrix},
\]
we obtain
\[
w_{k+1,1}
=
-\frac{
\eta v_{k+1}^R\,w_{k,1}
-(v_{k+1}^R\bar c+\eta\beta_{k+1}^R)w_{k,2}
}{
\Delta_{k+1}^R
}.
\]
Therefore
\[
\overline{w_{k,2}}\,w_{k+1,1}
=
-\frac{
\eta v_{k+1}^R\,\overline{w_{k,2}}\,w_{k,1}
-(v_{k+1}^R\bar c+\eta\beta_{k+1}^R)|w_{k,2}|^2
}{
\Delta_{k+1}^R
}.
\]
Taking imaginary parts and using
\[
\Im(\overline{w_{k,2}}\,w_{k,1})=-j_k,
\qquad
\Im(\bar c)=qy,
\qquad
\Im(\beta_{k+1}^R)=\tau_{k+1}^R,
\]
yields
\[
h_k
=
\frac{
\eta v_{k+1}^R\,j_k
+
(\eta\tau_{k+1}^R+qy\,v_{k+1}^R)|w_{k,2}|^2
}{
\Delta_{k+1}^R
}.
\]

Similarly, from the left-tail identity
\[
w_{k-1}=-(\Gamma_{k-1}^{-1}A)w_k,
\]
with
\[
\Gamma_{k-1}^{-1}
=
\frac1{\Delta_{k-1}^L}
\begin{pmatrix}
V_{k-1}^L & -\gamma_{k-1}^L\\
-\overline{\gamma_{k-1}^L} & U_{k-1}^L
\end{pmatrix},
\qquad
A=
\begin{pmatrix}
\eta & 0\\
-c & \eta
\end{pmatrix},
\]
we obtain
\[
w_{k-1,2}
=
-\frac{
(-\eta\overline{\gamma_{k-1}^L}-U_{k-1}^Lc)\,w_{k,1}
+\eta U_{k-1}^L\,w_{k,2}
}{
\Delta_{k-1}^L
}.
\]
Hence
\[
\overline{w_{k-1,2}}\,w_{k,1}
=
\frac{
(\eta\gamma_{k-1}^L+U_{k-1}^L\bar c)|w_{k,1}|^2
-\eta U_{k-1}^L\,\overline{w_{k,2}}\,w_{k,1}
}{
\Delta_{k-1}^L
}.
\]
Taking imaginary parts gives
\[
h_{k-1}
=
\frac{
(\eta\tau_{k-1}^L+qy\,U_{k-1}^L)|w_{k,1}|^2
+\eta U_{k-1}^L\,j_k
}{
\Delta_{k-1}^L
}.
\]
The positivity follows immediately.
\end{proof}

\begin{proposition}[Direct calculation of \(\Psi_k^\pm\)]
Let
\[
p_j\coloneqq -\Im(u_j u_{n+1-j}),
\qquad
s\coloneqq \sigma_{\min}(x+iy-A_n).
\]
Then for every interior block \(2\le k\le m-1\),
\[
s\,p_{2k-1}
=
y|w_{k,1}|^2+a h_{k-1}-b j_k,
\]
\[
s\,p_{2k}
=
y|w_{k,2}|^2+a j_k-b h_k.
\]
Consequently,
\[
s\,p_{2k-1}=|w_{k,1}|^2\,\Psi_k^-,
\qquad
s\,p_{2k}=|w_{k,2}|^2\,\Psi_k^+,
\]
where
\[
\Psi_k^-
=
y
+
\frac{a(\eta\tau_{k-1}^L+qy\,U_{k-1}^L)}{\Delta_{k-1}^L}
+
\left(
\frac{a\eta U_{k-1}^L}{\Delta_{k-1}^L}-b
\right)\frac{t_k}{D_k},
\]
and
\[
\Psi_k^+
=
y
-
\frac{b(\eta\tau_{k+1}^R+qy\,v_{k+1}^R)}{\Delta_{k+1}^R}
+
\left(
a-\frac{b\eta v_{k+1}^R}{\Delta_{k+1}^R}
\right)\frac{t_k}{A_k}.
\]
Equivalently,
\[
p_{2k-1}>0 \iff \Psi_k^- >0,
\qquad
p_{2k}>0 \iff \Psi_k^+ >0.
\]
\end{proposition}

\begin{proof}
The Takagi site identities are
\[
s\,p_j
=
y|u_j|^2+a I_{j-1}-b I_j.
\]
For \(j=2k-1\), this becomes
\[
s\,p_{2k-1}
=
y|w_{k,1}|^2+a h_{k-1}-b j_k.
\]
For \(j=2k\), it becomes
\[
s\,p_{2k}
=
y|w_{k,2}|^2+a j_k-b h_k.
\]
Now substitute the formulas for \(h_{k-1}\) and \(h_k\) from the previous proposition, and use
\[
j_k=\frac{t_k}{D_k}|w_{k,1}|^2=\frac{t_k}{A_k}|w_{k,2}|^2.
\]
After factoring out \(|w_{k,1}|^2\) and \(|w_{k,2}|^2\), one obtains the displayed formulas for
\(\Psi_k^\pm\).
\end{proof}

\begin{corollary}[Even-\(n\) symmetry reduction]
If \(n=2m\), then
\[
p_j=p_{2m+1-j}
\qquad (1\le j\le 2m).
\]
Hence
\[
p_{2k-1}=p_{2(m-k+1)},
\qquad
p_{2k}=p_{2(m-k)+1},
\]
so it is enough to prove positivity of one branch, for example
\[
\Psi_k^+>0
\qquad (1\le k\le m).
\]
\end{corollary}

\begin{proof}
This is immediate from
\[
p_j=-\Im(u_j u_{2m+1-j}).
\]
\end{proof}

\end{document}
